\inicializarSeccion{Introducción}
\tab En este documento se desarrolló el reto de la aplicación de campos eléctricos al diagnóstico de la malaria mediante la dielectroforesis. El objetivo principal de este proyecto es analizar el comportamiento de los campos eléctricos como herramientas para la manipulación de la biomolécula y en sus posibles aplicaciones del diagnóstico de esta enfermedad.

\tab Al principio se investigó la interacción de los glóbulos rojos y de los campos eléctricos, así como el impacto de la malaria en la población mundial. Después se desarrolló una modelación computacional en el software MATLAB para simular el campo eléctrico generado por configuraciones de las cargas, comenzando con un dipolo y extendiéndose la distribución en las líneas de cargas.

\tab Finalmente, se implementa la interfaz gráfica utilizando el AppDesigner, lo que permite al usuario poder modificar los parámetros de la interfaz gráfica, también poder visualizar el campo eléctrico en términos de magnitud y dirección, y analizar su comportamiento de forma interactiva. Esto con el fin de facilitar la comprensión de los campos eléctricos no uniformes fundamentales en la aplicación de la dielectroforesis.
